\section{Related Work}
\label{sec:related-work}
% ================================================================================

This section surveys the protocols and implementations closest to the
Bridge Contract and states what each leaves uncovered.

\subsection{Agent Transfer Protocol (AGTP)}\label{sec:rw-agtp}

AGTP \cite{I-D.hood-agtp} is an application-layer protocol for AI agent
traffic that introduces agent-native methods (among them QUERY, SUMMARIZE, EXECUTE,
DELEGATE, CONFIRM, ESCALATE, SUSPEND, NOTIFY) and protocol-level identity
headers including \texttt{Agent-ID}, \texttt{Owner-ID}, and
\texttt{Authority-Scope}; \texttt{Session-ID} is an optional header.
\texttt{Owner-ID} is the principal identifier carried in the agent
identity document referenced by \texttt{Agent-ID} and surfaced as the
\texttt{Owner-ID} request header when explicit wire-layer carriage is
needed (AGTP-09 \S 2, Terminology); \texttt{Agent-ID} and \texttt{Authority-Scope}
are also request headers (AGTP-09 \S 6). Its
session companion \cite{I-D.hood-agtp-session} distinguishes bounded and
persistent sessions (\S 3) and inherits identity, authority, and
attribution from base AGTP (\S 5).

AGTP defines a pause and resume pair at method level. SUSPEND
(AGTP-09 \S 8.2.12) places a session workflow in a recoverable paused
state and returns a single-use resumption nonce of at least 128 bits,
an optional \texttt{resume\_by} deadline after which the session
expires, and an optional agent-provided checkpoint that the server
stores while the session is suspended. It is distinct from the
lifecycle SUSPEND of the agent's registry record. The corresponding
RESUME method is named in AGTP-09 \S 8.2.12 (the SUSPEND response
notes that ``once presented to resume a session, the nonce is
invalidated''), but AGTP-09 does not specify the RESUME request
message itself. AGTP-09 lists RESUME among its Tier-2 ORCHESTRATE
methods (\S 8.6.3). The companion vocabulary
\cite{I-D.hood-agtp-standard-methods}, a separate draft from the AGTP
base that could not be retrieved for this paper, is reported to list
RESUME in its \S 8.7 (ORCHESTRATE category); we did not read it. The Bridge Contract concerns the
case this pair does not: an abrupt detach in which no SUSPEND was
issued. As read, AGTP-09 SUSPEND binds the resumption nonce to no
audit position (the nonce is returned in a method response and is not
tied to a chain sequence or a governance event), and AGTP and its
companions define no handling of tool calls or human confirmations
that arrive during the pause and no relation between a detach timeout
and the human wait. AGTP positions itself as a ``narrow waist''---a
transport-layer construct that other agent frameworks compose with
rather than replace. AGTP does not specify what a bridge component
sitting between a client transport and an agent runtime must do when the
client transport disconnects while the agent continues to execute, nor
when the human in the loop is unavailable. The Bridge Contract defines
the execution semantics for AGTP headers in exactly those scenarios; it
does not redefine the headers themselves.

\paragraph{Why this obligation is out of AGTP's scope.}
AGTP-09 is a transport-layer ``narrow waist''; bridge execution
during detach is a governance obligation at a layer above it. A
future revision that added bridge semantics would be expanding its
scope, not filling a gap inside it.

\subsection{Governance Audit Record (GAR)}\label{sec:rw-gar}

GAR \cite{I-D.sato-soos-gar} defines a tamper-evident audit architecture
comprising its audit types, the Session Audit Record (SAR), the Audit
Alert system, and the Audit Package. GAR assumes a Governing Enforcement
Component (GEC) that is the sole source of audit truth and that signs
its audit artifacts, among them Session Audit Records and the Merkle
roots of Session Blocks.
GAR's design premise is non-suppressibility: the GEC ``MUST generate audit
artifacts automatically, MUST sign them, and MUST NOT allow any agent,
application, or principal to suppress, modify, or delete them''
(GAR-08 \S 1, Introduction). The premise is sound. We
introduce the following deployment profile to clarify the Bridge
Contract's positioning with respect to GAR.

\begin{definition}[GEC-less deployment profile]\label{def:gecless}
A \emph{GEC-less deployment profile} is a deployment in which no
Governing Enforcement Component (GEC) is present and in which the
bridge is the sole emitter of governance events for sessions it
mediates. In a GEC-less deployment the bridge's event log is the
authoritative source of session governance events for downstream
consumers (e.g., a SCITT transparency log or an external auditor).
The bridge's events are emitted in a \emph{GAR-shaped} form---they
are causally ordered, hash-chained, and serialised as canonical JSON
(JSON Canonicalization Scheme, JCS, RFC~8785~\cite{rundgren2020json})---but
are not signed by a GEC. The chain construction follows the Agent Audit
Trail's \texttt{prev\_hash} over RFC~8785 (\S\ref{sec:rw-audit-trail}),
not GAR's \texttt{prev\_span\_hash}. Event signing in a GEC-less deployment is governed
by Invariant I-SIGN-EVENT (Section~\ref{sec:invariants-signing}).
\end{definition}

The Bridge Contract does \emph{not} claim to be a substitute for a
GEC. In GEC-present deployments the bridge forwards events to the GEC
for signing; the GEC remains the sole source of audit truth. In
GEC-less deployments the bridge is the emitter, but consumers of its
events MUST verify the operator key binding out-of-band (this paper's
rule, by analogy with GAR \S 6.1). The Bridge Contract is thus a \emph{GAR-shaped event
source for GEC-less deployments}, not a replacement for GAR's GEC-based
audit architecture.

\paragraph{Why this obligation is out of GAR's scope.}
GAR-08 premises its architecture on a GEC; the GEC-less emitter is
a deployment profile, not a GAR concern. A future revision that
defined a GEC-less mode would be lowering its own premise, not
filling a gap inside its scope.

\subsection{Human Escalation Mechanism (HEM)}\label{sec:rw-hem}

HEM \cite{I-D.sato-soos-hem} is a companion of GAR by the same author and
defines the human side of a governed session. A trigger moves the session
to \texttt{HEM\_PENDING} (\S 4, \S 5), during which the GEC MUST reject
every state transition (\S 13.1). Each principal in a designation chain has
a timeout budget with a minimum of 60 seconds and a pre-declared
disposition: escalate to the next principal, suspend, terminate, or, under
Cedar-gated conditions, auto-approve (\S 14.1, \S 14.2). A decision is a
submission naming the \texttt{hem\_id}, signed with Ed25519, and admitted
at most once; a later submission against a resolved event is rejected, and
the exact action is checked by verifying the decision and the
kernel-signed request jointly (\S 12.8). This overlaps the contract's
single-use answers and its timer rules, and the contract does not present
those mechanisms as new. HEM's absent party is the human principal, reached
through a delivery channel, not a client transport. As read, it defines no
queue of tool calls held while a session's client is detached, no resume
record or re-attach proof, and no relation between a detach timeout and
the human wait. It is written against a GEC kernel, which the bridge of
the contract is not.

\paragraph{Why this obligation is out of HEM's scope.}
HEM-07 is the human side of a GAR-governed session; the client side
and a non-GEC bridge component are out of its scope. A HEM profile
that owned them would have to introduce a non-GEC component and a
client-transport dimension; that is a scope expansion.

\subsection{Protocol for Agent Communication Exchange (PACE)}\label{sec:rw-pace}

PACE \cite{I-D.nandakumar-agentproto-moq-pace} defines a session protocol
for AI agent communication over Media over QUIC Transport (MOQT). Its
session recovery (\S 4.5) resumes a SUSPENDED session when the agent
reconnects to the relay and re-subscribes on its agent prefix, after which
the relay delivers cached objects from the point of disconnection; agents
SHOULD publish state checkpoints. PACE's message-type registry (its
Table~17) reserves \texttt{SESSION\_SUSPEND} (0x05) and
\texttt{SESSION\_RESUME} (0x06), but, as read, defines no message body
for either and no resume credential that a client must present; the
\texttt{SESSION\_RESUME} row reads only ``Resume suspended session''.
PACE also supports delegated authorization with scope narrowing on each
delegation (\S 9.1.2); the related \S 5.1.2 concerns delegation-chain
identity tracking, not scope narrowing. PACE presumes a relay or
MOQT substrate that caches objects while the client is gone. PACE does not specify what happens during the detach window when
the agent continues to execute and produces additional state that must be
reconciled on re-attach. The Bridge Contract's Resume Record (Definition~\ref{def:rt}) is the
bridge's own causally-pinned snapshot of a detachment, and its re-attach
protocol (challenge-response over the last event the client saw)
defines what a bridge checks before it replays; a bridge MAY export the
record to a PACE-compatible relay for cross-vendor or cross-relay
interoperation, as a non-authoritative snapshot; that export is optional and not
specified here, so the cross-vendor case is aspirational until a second
implementation exists. Intra-PACE recovery---where the same relay
holds the suspended session throughout---does not require the Bridge
Contract's resume record; the contract is necessary only
when the suspension window spans relay restarts, vendor handoff, or
cross-relay re-attach. The Bridge Contract thus complements PACE: PACE
describes relay-level reconnection (\S 4.5); the Bridge Contract specifies ``what
happened during the detach'' for the cross-vendor case.

\paragraph{Why this obligation is out of PACE's scope.}
PACE-00 is resumption over MOQT; bridge governance during absence is
a different layer. A PACE profile that owned the Bridge Contract's
obligations would have to specify a non-relay bridge component, a
governance event log, and a challenge-response proof tied to that
log; that is a different architecture from MOQT relaying.

\subsection{Agent Audit Trail}\label{sec:rw-audit-trail}

The Agent Audit Trail \cite{I-D.sharif-agent-audit-trail} defines a
record format with registered action types (among them \texttt{tool\_call},
\texttt{decision}, \texttt{delegation}, \texttt{escalation} and
\texttt{lifecycle}), outcomes (success, failure, timeout, denied,
escalated), trust levels L0--L4, and a chain in which \texttt{prev\_hash}
is the SHA-256 of the canonical JSON (RFC~8785) of the previous record
(\S 6, \S 8.2). It is session-structured: a session begins with a genesis
record and ends with a close record carrying a session hash (\S 8.1,
\S 8.3); lifecycle records can log pause and resume (\S 7); a record can be
written before the action it authorises (\S 4); and a crash is documented by
a recovery record or a synthetic close for an orphaned session
(\S 8.2, \S 8.3). What it specifies is the record and the chain. As read, it
defines no proof that a re-attaching client saw a given prefix of the chain,
no queue of calls awaiting a human while the client is away, and no timers
for that wait. The Bridge Contract's events chain the same way
(\texttt{prev\_hash} over canonical JSON) but carry a per-session
\texttt{sequence} number rather than a \texttt{parent\_record\_id}; no
mapping between the two schemas is specified here. The contract adds the
session-level state machine, the resume record and the re-attach protocol.

\paragraph{Why this obligation is out of the Agent Audit Trail's scope.}
The Agent Audit Trail is a record format and chain; bridge execution
and re-attach proof are session-protocol concerns at a different
layer.

\subsection{Agent-to-Agent (A2A) and Model Context Protocol (MCP)}

A2A defines a task lifecycle that includes an \texttt{input-required}
state (not pinned in the source manifest), which overlaps the Bridge Contract's \texttt{HUMAN\_PENDING}
cell. MCP's Streamable HTTP transport defines session identifiers,
session termination, and \texttt{Last-Event-ID} resumability (not
pinned in the source manifest), which
overlaps the Bridge Contract's re-attach mechanism. Neither A2A nor
MCP specifies what a bridge component MUST do during client detach,
including durable event logging, queue management, timer-driven
termination, or challenge-response re-attach pinned to an audit
chain. The Bridge Contract complements both: A2A and MCP define the
messaging and task lifecycle; the Bridge Contract defines the
session-governance obligations of an intermediary that holds them
during detach.

\subsection{Herdr}

Herdr \cite{herdr-product} is a product (Apache-2.0, Rust-based) that
keeps terminals running in a background server when the client
disconnects, restores the saved layout and can resume supported agent
sessions after a restart, offers an experimental live handoff, and keeps
layout snapshots. It is an implementation, not a specification. Its
documentation (version~0.9.5) names \texttt{session.json} and the
snapshot files but, as read, gives no schema, compatibility promise or
conformance criteria for them, and snapshots hold layout and agent session
references, not running processes or history. We use it as
evidence of demand for the bridge responsibility class; the relationship
is one of specification to implementation, not competition.

\subsection{AgentThread}\label{sec:rw-agtp-thread}

AgentThread \cite{zheng2026formal} (Zheng et al., arXiv:2606.28690) is a
source-linked framework for security assurance analysis of agent
protocols that compiles specification text into TLA+-style checks and
replays counterexamples against SDKs. Its central finding is that no
protocol assigns enforcement for cross-protocol behaviour, a
responsibility gap across protocols, SDKs, and deployments. This
motivates the Bridge Contract; it is not a result about AGTP, GAR or PACE.
AgentThread analyses the protocols MCP, A2A, ANP, and two different
protocols both abbreviated ``ACP'' in the source paper --- the
\emph{Agent Communication Protocol} (ACP) and the \emph{Agent Client
Protocol} (ACP-Client). Throughout this document, ``ACP'' means the
\emph{Agent Client Protocol} (the protocol whose
\texttt{session/request\_permission} the bridge may see, see
Section~\ref{sec:coverage}); the AgentThread paper itself uses
``ACP-Client'' for that protocol and ``ACP'' for the Agent Communication
Protocol. AgentThread provides the
analysis but no specification; the Bridge Contract specifies one class of
the orphaned obligations, those at the bridge layer.

\subsection{Precedents}

Several techniques used by the Bridge Contract have well-established
precedents that this document does not claim as novel:

\begin{itemize}[leftmargin=2em]
\item \textbf{Ack-and-replay resumption.} XMPP stream management
(XEP-0198) and sequence-numbered gateway RESUME protocols use
acknowledgement and replay to recover from transport disconnect.
The Bridge Contract's re-attach replay (I4) is in the same family,
but adds challenge-response authentication and hash-chain prefix
verification, which XEP-0198 does not.
\item \textbf{Tamper-evident logging.} Crosby and Wallach's
work on append-only logs, RFC~6962 (Certificate Transparency), and
RFC~9162 (Certificate Transparency v2) define Merkle-tree-based
tamper-evidence. The Bridge Contract's I5 hash chain and
I-CHECKPOINT Merkle root follow the same construction; the
contribution is the integration with session governance, not the
Merkle tree itself.
\item \textbf{Durable execution and checkpoint-and-interrupt.}
Workflow engines and durable-execution frameworks checkpoint state
and resume after interruptions. The Bridge Contract's resume record
and restart recovery are in the same family, but add the
session-governance invariants (authority lock, queue integrity,
answer binding) that workflow engines do not specify.
\item \textbf{AgentRFC.} AgentRFC (arXiv:2603.23801) is named here
as related prior art on agent protocol composition; it was not read
for this paper and no claim is made about its contents.
\end{itemize}

\subsection{Non-duplication statement}

Each protocol above leaves a bridge obligation uncovered: AGTP has
transport headers but no bridge execution during detach; GAR assumes a
GEC as sole event source (see Definition~\ref{def:gecless}); PACE describes relay-level
reconnection (\S 4.5) but not the detach window; HEM specifies human escalation for a
governed object but not a detached client; A2A and MCP operate above the
session layer. Herdr is a product and AgentThread an analysis, neither a
specification. The ``Why this obligation is out of <draft>'s scope''
paragraphs in \S\ref{sec:rw-agtp}, \S\ref{sec:rw-gar},
\S\ref{sec:rw-hem}, \S\ref{sec:rw-pace} and the Agent Audit Trail
subsection argue that this is structural (a scope mismatch), not
historical (a gap that a future revision might fill): a future
revision of any of these drafts that owned the Bridge Contract's
obligations would be expanding its scope, not closing a gap inside
it. That argument is what makes the non-duplication claim robust to
future revisions of the cited drafts.

\paragraph{Comparison table.}
Table~\ref{tab:comparison} states, for each cited source, what the Bridge
Contract adds and whether that addition is necessary (no other draft
provides it) or a design choice (the draft could have been extended to
provide it). What each source requires and leaves open is in the
subsections above, and the obligation-by-obligation comparison, with its
method, is the matrix in Section~\ref{sec:mapping}. Every entry is pinned to
the version listed in \texttt{paper/sources/MANIFEST.md}.

\begin{table}[htbp]
\centering
\caption{What the Bridge Contract adds to each cited source, and whether
the addition is necessary or a design choice.}
\label{tab:comparison}
\footnotesize
\begin{tabularx}{\columnwidth}{l X X}
\toprule
\textbf{Draft} & \textbf{Bridge Contract adds} & \textbf{Necessary?} \\
\midrule
AGTP-09 &
A bridge-side state machine, the resume record bound to a chain
position (I1, I4), and the challenge-response re-attach that uses it. &
Necessary (scope, §\ref{sec:rw-agtp}): AGTP-09 is a transport narrow waist; bridge execution during detach is a layer above it. \\
AGTP standard-methods-01 (\S 8.7) &
The challenge-response re-attach that does not require a prior SUSPEND,
and that pins the prefix the client last saw. &
Necessary (scope, §\ref{sec:rw-agtp}): the companion draft is reported to list a RESUME method for a suspended session (not read); the contract resumes one that was detached without a SUSPEND, which is a different layer. \\
AGTP-SESSION-00 &
$T_{\text{detach}}$, $T_{\text{total}}$, $T_{\text{human}}$ and the
detach/re-attach state machine. &
Necessary (scope, §\ref{sec:rw-agtp}): AGTP-SESSION governs session establishment, not bridge execution during the detach window. \\
AGTP-BINDINGS-00 &
Invariant I8, which makes the bridge reject any re-attach the transport
reports as a replay. &
Design choice: the binding could in principle extend to specify the
bridge-layer consequence; the contract states it directly. \\
GAR-08 &
The GEC-less deployment profile (Definition~\ref{def:gecless}) and the
bridge as the event source for that case, emitting in GAR shape. &
Necessary (scope, §\ref{sec:rw-gar}): GAR is an audit architecture premised on a GEC; the GEC-less emitter is a deployment profile, not a GAR concern. \\
HEM-07 &
The tool-call queue (I2, I13), the answer binding (I14), and the
challenge-response re-attach. &
Necessary (scope, §\ref{sec:rw-hem}): HEM is the human side of a GAR-governed session; the client side and a non-GEC bridge component are out of HEM's scope. \\
PACE-00 &
The resume record $R$ (Definition~\ref{def:rt}) and the
challenge-response re-attach that pins the prefix the client last saw. &
Necessary (scope, §\ref{sec:rw-pace}): PACE is resumption over MOQT; bridge governance during absence is a different layer. \\
Agent Audit Trail-01 &
The challenge-response re-attach pinned to the prefix (I4) and the
tool-call queue with its timers (I2, I13). &
Necessary (scope, §\ref{sec:rw-audit-trail}): the audit trail is a record format and chain; bridge execution and re-attach proof are session-protocol concerns. \\
Herdr 0.9.5 &
A specification with conformance criteria. &
Design choice: a product could publish its format as a de-facto
standard; the contract chooses to specify the bridge obligations
abstractly. \\
AgentThread (arXiv:2606.28690) &
A specification of one class of those obligations, at the bridge layer. &
Necessary: AgentThread is an analysis framework, not a specification. \\
\bottomrule
\end{tabularx}
\end{table}

The table reads each cited draft at the version listed in
\texttt{paper/sources/MANIFEST.md}. The ``Necessary'' column marks an
addition that no cited draft provides; ``Design choice'' marks one that
could in principle have been added to an existing draft but is stated
here for clarity.
